\section{失败案例与调试手册：把不可用系统变成可迭代系统}

\subsection{失败类型分层：先分桶再修复}
在多模态 Agent 项目里，失败不是例外而是常态。课程隐含的一个关键工程思想是：不要直接盯住 ``最终失败''，而要先把失败拆成可修复子类。只有当失败类型被分桶，优化动作才会高效。

\begin{importantbox}{建议的四层失败分类}
\begin{enumerate}
    \item \textbf{感知失败}：没看懂界面、没读到关键文本、目标对象定位错误；
    \item \textbf{规划失败}：子目标顺序错、漏步骤、重复步骤过多；
    \item \textbf{执行失败}：动作格式对但上下文不匹配，或接口调用失败；
    \item \textbf{验证失败}：任务已接近完成，但校验脚本与目标定义不一致。
\end{enumerate}
\end{importantbox}

\begin{knowledgebox}{为什么 ``失败分桶'' 会直接影响训练效率}
如果所有失败样本都被统一回流，会造成噪声放大。分桶后可以按问题类型定向补数据：感知失败补 grounding 与 OCR 样本，规划失败补长链路推理样本，执行失败补动作约束样本，验证失败补脚本鲁棒性测试样本。
\end{knowledgebox}

\subsection{调试闭环：日志、回放、对照实验}
课程内容对应的工程实践可以抽象为三步：第一步抓全量日志，第二步做轨迹回放，第三步做最小对照实验。对照实验的目标不是追求最佳分数，而是确认某个改动是否真实减少了某类失败。

\begin{table}[H]
\centering
\begin{tabular}{p{2.6cm}p{4.6cm}p{4.0cm}}
\toprule
\textbf{阶段} & \textbf{关键操作} & \textbf{输出物} \\
\midrule
日志采集 & 保存观察、动作、环境状态、评测脚本输出 & 可定位失败发生点 \\
轨迹回放 & 按时间序列回放 Agent 行为与界面变化 & 失败原因初步归因 \\
对照实验 & 单变量替换模型/提示/动作约束/脚本版本 & 可验证的修复证据 \\
\bottomrule
\end{tabular}
\caption{多模态 Agent 的最小可行调试流程}
\end{table}

\begin{warningbox}{避免 ``一次改很多''}
当模型、数据、动作格式、评估脚本同时变化时，几乎无法判断真实增益来源。课程强调工程化迭代，本质上要求每轮实验保持可解释和可回退。
\end{warningbox}

\subsection{高频问题实例}
\begin{itemize}
    \item \textbf{元素定位偏移}：界面缩放变化导致坐标偏差，建议引入相对定位与对象锚点；
    \item \textbf{长任务退化}：前半段正确、后半段崩溃，常见于上下文压缩和规划遗忘；
    \item \textbf{奖励欺骗}：Agent 学会触发脚本漏洞而非完成任务，需要增强验证脚本鲁棒性；
    \item \textbf{跨应用状态错位}：多个窗口与后台任务并行导致状态不一致，需要更强会话管理。
\end{itemize}

\begin{importantbox}{调试优先级排序}
先修复 ``高频 + 高损失'' 类失败，再处理低频边缘问题。对生产系统而言，稳定性收益通常大于单次最优表现。
\end{importantbox}

\subsection{本章小结}
失败调试不是附属工作，而是 Agent 迭代主流程。把失败结构化、可回放、可对照，才能形成真正的训练飞轮。

